\documentclass[12pt,a4paper]{article}
\usepackage[utf8]{inputenc}
\usepackage[T1]{fontenc}
\usepackage[ngerman]{babel}
\usepackage{amsmath,amssymb}
\usepackage{geometry}
\geometry{a4paper,left=2cm,right=2cm,top=2cm,bottom=2cm}
\usepackage{parskip}
\usepackage{multicol}
\usepackage{array}

\begin{document}

% === FRONT PAGE (VORDERSEITE) ===
\begin{center}
    \Large\textbf{WEBER-DE BROGLIE-BOHM-THEORIE (WDBT+)} \\
    \large{Ein alternatives physikalisches Paradigma}
\end{center}
\vspace{0.5cm}

\section*{Kernaussagen}
\begin{itemize}
    \item \textbf{Feldfrei:} Keine elektromagnetischen Felder, keine Raumzeitkrümmung -- nur direkte Teilchenwechselwirkung
    \item \textbf{Deterministisch:} Kein Zufall, keine Wellenfunktionskollapse -- Führungswelle mit Quantenpotential
    \item \textbf{Singularitätsfrei:} Quantenpotential verhindert unendliche Dichten (Schwarze Löcher, Urknall)
    \item \textbf{Fraktaler Raum:} Fundamentale Dimension $D \approx 2{,}71$ (ermöglicht Ableitung der Naturkonstanten)
    \item \textbf{Emergenz:} Alle bekannten Theorien sind Grenzfälle der WDBT+
\end{itemize}

\section*{Grundgleichungen}

\subsection*{1. Weber-Elektrodynamik (WED) -- Erweiterung der Coulomb-Kraft}
\begin{equation}
F_{\text{WED}} = \frac{q_1q_2}{4\pi\epsilon_0 r^2}\left[1 - \frac{\dot{r}^2}{c^2} + 2\frac{r\ddot{r}}{c^2}\right]
\end{equation}
($\beta = 2$ ist fest in der Elektrodynamik)

\subsection*{2. Weber-Gravitation (WG) -- Erweiterung des Newton-Gesetzes}
\begin{equation}
F_{\text{WG}} = -\frac{GMm}{r^2}\left[1 - \frac{\dot{r}^2}{c^2} + \beta\frac{r\ddot{r}}{c^2}\right] \hat{r}
\end{equation}
\begin{itemize}
    \item \textbf{Reine WG:} $\beta = 1$ für Photonen (Lichtablenkung), $\beta = 0{,}5$ für Massen
    \item \textbf{Mit DBT:} $\beta = 0{,}5$ für \textit{alle} Teilchen (Quantenpotential übernimmt restliche Wirkung)
\end{itemize}

\subsection*{3. Quantenpotential (DBT)}
\begin{equation}
Q = -\frac{\hbar^2}{2m}\frac{\nabla^2R}{R}
\end{equation}
Führt Teilchen deterministisch; ergibt mit WED/WG die vollständige Bewegung:
\begin{equation}
m\ddot{x} = -\nabla V - \nabla Q
\end{equation}

\subsection*{4. Fraktale Raumdimension}
\begin{equation}
D = \frac{\ln 20}{\ln(2+\phi)} \approx 2{,}71, \qquad \phi = \frac{1+\sqrt{5}}{2}
\end{equation}
Fundamentale Struktur, aus der Naturkonstanten folgen.

\section*{Was emergiert}
\begin{itemize}
    \item \textbf{Maxwell-Gleichungen \& Lorentz-Kraft} (Mittelung der WED)
    \item \textbf{Spezielle Relativität} (aus WG + kosmischer Mittelung)
    \item \textbf{Schrödinger-Gleichung} (Madelung-Transformation der DBT)
    \item \textbf{ART} (durch Quantenpotential singularitätenfrei erweitert)
\end{itemize}

\section*{Was wird erklärt (ohne Dunkle Materie)}
\begin{itemize}
    \item Galaxienrotation (WG + fraktale Struktur)
    \item CMB-Anomalien (fraktale Dimension)
    \item Plasma-Instabilitäten (nicht-lokale WED)
    \item Actio = Reactio (WED erfüllt Impulserhaltung, Maxwell nicht)
\end{itemize}

\newpage
% === BACK PAGE (RÜCKSEITE) ===
\begin{center}
    \Large\textbf{WEBER-DE BROGLIE-BOHM-THEORIE (WDBT+)} \\
    \large{Erweiterungen, Vorhersagen \& kosmologisches Paradigma}
\end{center}
\vspace{0.5cm}

\section*{Dynamische Trägheits-Schwere Theorie (DSTT)}
Die DSTT ist eine direkte Konsequenz der Weber-Struktur und des Quantenpotentials. Sie besagt, dass sich die Trägheit in eine radiale und eine tangentiale Komponente aufspaltet:

\begin{equation}
m\ddot{\vec{r}} = -\nabla V_{\text{radial}} - \nabla V_{\text{tangential}}
\end{equation}

\textbf{Konsequenzen:}
\begin{itemize}
    \item \textbf{Keine geschlossenen Bahnen:} Alle Umlaufbahnen sind Spiralen (nach innen oder außen)
    \item \textbf{Erklärung der Planetenbewegung:} Merkurs Periheldrehung, Jupiterwanderung, Hot Jupiters
    \item \textbf{Galaxienrotation:} Flache Rotationskurven ohne Dunkle Materie
\end{itemize}

\section*{Gravitation als Maxwell-Theorie}
Aus der DSTT und der Weber-Gravitation folgt eine \textbf{Maxwell-Form der Gravitation} mit zwei Komponenten:
\begin{itemize}
    \item Radiale Komponente $\leftrightarrow$ elektrisches Feld
    \item Tangentiale Komponente (Spiralform) $\leftrightarrow$ magnetisches Feld
\end{itemize}
Damit wird die strukturelle Einheit von Elektrodynamik und Gravitation sichtbar – beide gehorchen derselben mathematischen Form.

\section*{Entropie der Gravitation}
Da alle Bahnen spiralisieren (irreversibel), produziert Gravitation permanent Entropie:
\begin{itemize}
    \item Keine ewigen Planetenbahnen – alles hat eine endliche Lebensdauer
    \item Gravitation ist nicht zeitumkehrsymmetrisch
    \item Die Entropieproduktion treibt den kosmischen Energiezyklus an
\end{itemize}

\section*{Kosmologisches Paradigma: Steady State \& Energiezyklus}
Die WDBT+ führt zu einem geschlossenen, zyklischen Universum ohne Urknall:

\begin{itemize}
    \item \textbf{Keine Expansion:} Die Rotverschiebung ist nicht Doppler-Effekt, sondern kumulative Gravitation (WG)
    \item \textbf{Energiekreislauf:} 
    \[
    \text{Photonen} \xrightarrow{\text{Rotverschiebung}} \text{Nullpunktenergie} \xrightarrow{\text{Materieentstehung}} \text{Neue Photonen}
    \]
    \item \textbf{Kontinuierliche Materieentstehung} aus dem Quantenpotential des Vakuums
    \item \textbf{Erklärung der JWST-Daten:} Galaxien bei $z \approx 25$ sind nicht "zu früh" – sie sind alt, die Rotverschiebung täuscht keine Jugend vor
    \item \textbf{Achse des Bösen (CMB-Anomalie):} Natürliche Folge der fraktalen Dimension $D \approx 2{,}71$
\end{itemize}

\section*{Entscheidende Experimente}

\subsection*{1. Frequenzabhängige Lichtablenkung an Massen}
\begin{itemize}
    \item ART: Rein geometrisch $\rightarrow$ Ablenkungswinkel \textbf{frequenzunabhängig}
    \item WDBT+: Dynamisch durch WG $\rightarrow$ \textbf{$\Delta\varphi \propto \lambda^2$} (messbar bei Radiogalaxien)
\end{itemize}

\subsection*{2. Modifizierter Lamb-Shift}
\begin{itemize}
    \item QED: Bisherige Präzisionsrechnungen
    \item WDBT+: Zusatzterm durch Vakuum-Quantenpotential:
    \begin{equation}
    \Delta E_{\text{Lamb}}^{\text{WDBT}} = \Delta E_{\text{QED}} + \frac{e^2\hbar}{4\pi\epsilon_0 m_e^2c^3}\langle r\rangle
    \end{equation}
\end{itemize}

\subsection*{3. Shapiro-Laufzeitverzögerung}
\begin{itemize}
    \item ART: Laufzeit unabhängig von Signalfrequenz
    \item WDBT+: \textbf{Frequenzabhängige Verzögerung} (testbar mit Pulsarsignalen)
\end{itemize}

\subsection*{4. Hubble-Gesetz auf großen Skalen}
\begin{itemize}
    \item Standard: Linear (Expansion des Raums)
    \item WDBT+: \textbf{Abweichungen vom linearen Verlauf} (kumulative Gravitation, keine Expansion)
\end{itemize}

\subsection*{5. Fraktale Signaturen im Plasma}
\begin{itemize}
    \item Vorhersage: Stromdichten skalieren mit \textbf{$j(r) \propto r^{-0{,}29}$} (aus $D \approx 2{,}71$)
    \item Nachweisbar in: Laborplasmen, Sonnenwind, Galaxienfilamenten
\end{itemize}

\subsection*{6. Spiralbahnen statt Ellipsen}
\begin{itemize}
    \item Präzise Bahnverfolgung von Satelliten und Planeten müsste langsame Spiralbewegung zeigen (radiale Drift)
    \item Erklärung der Langzeitentwicklung von Planetensystemen
\end{itemize}

\section*{Vergleich mit Standardtheorien}
\begin{center}
\begin{tabular}{|p{3.0cm}|p{4.2cm}|p{4.5cm}|}
    \hline
    \textbf{Phänomen} & \textbf{ART / QM / Maxwell} & \textbf{WDBT+} \\
    \hline
    Gravitationswellen & Raumzeit-Vibrationen & Emergent aus fraktalem Raummodell, Spiralmoden \\
    \hline
    Schwarze Löcher & Singularität & Quantenpotential verhindert Kollaps $\rightarrow$ deterministisch \\
    \hline
    Dunkle Materie & Benötigt & Erklärt durch WG + fraktale Struktur + DSTT \\
    \hline
    Messproblem & Kollaps (unphysikalisch) & Deterministische Führungswelle (emergente QM) \\
    \hline
    Actio = Reactio & Maxwell verletzt es & WED erfüllt es exakt \\
    \hline
    Naturkonstanten & Empirisch (just so) & Folgen aus fraktaler Dimension $D \approx 2{,}71$ \\
    \hline
    Raumzeit & Fundamentale Krümmung & Emergent aus Teilchenwechselwirkung \\
    \hline
    Kosmologie & Urknall, Expansion & Steady State, Energiezyklus \\
    \hline
    CMB-Anomalien & Problem (benötigt Inflation) & Erklärt durch fraktale Dimension \\
    \hline
    JWST-Daten & Überraschung (zu frühe Galaxien) & Erwartet (keine Expansion) \\
    \hline
\end{tabular}
\end{center}

\section*{Konsequenzen}
\begin{itemize}
    \item \textbf{Emergenz:} ART, QM, Maxwell sind Grenzfälle einer tieferen Theorie
    \item Keine Raumzeit-Krümmung als Fundament (nur emergent)
    \item Keine Expansion des Universums (nur scheinbar)
    \item Keine Dunkle Materie (erklärt)
    \item Keine Singularitäten (verhindert)
    \item \textbf{Deterministische, nicht-lokale, feldlose Physik als Fundament}
    \item Gravitation hat Entropie – alle Bewegungen sind irreversibel
    \item Die Maxwell-Form der Gravitation schließt den Kreis zur Elektrodynamik
\end{itemize}

\vspace{0.3cm}
\noindent WDBT+ beansprucht: Alle etablierten Theorien emergieren aus einem gemeinsamen Rahmen. Die ART ist nicht falsch, sondern ein Grenzfall – so wie die Newtonsche Mechanik ein Grenzfall der ART ist. Die DSTT, der Energiezyklus und die fraktale Dimension erweitern das Bild zu einem vollständigen, in sich geschlossenen kosmologischen Paradigma.

\end{document}